% !TEX root = ../main.tex

A jelen fejezet a dolgozat empirikus részének alapját képező \textbf{YJMob100K} \cite{yabe2023metropolitan} adathalmazt mutatja be: az adatfelvétel körülményeit, a sémát, a két részadatkészlet közötti kapcsolatot, valamint a kísérő POI-fájlok szerepét. Az általános, kutatási célú CDR-adatkészletekre vonatkozó módszertani áttekintést a 3.7.\ szakasz tartalmazza.

\subsection{Az YJMob100K adatkészlet áttekintése}

Az YJMob100K adatkészletet Yabe és munkatársai \cite{yabe2023metropolitan} tették közzé 2023-ban a \emph{Scientific Data} folyóiratban, a HuMob Challenge mobilitás-előrejelzési verseny adatforrásaként. Az adatok a Yahoo Japan~Corporation egyik okostelefon-alkalmazásából (Yahoo Japan Application) származnak: anonimizált felhasználók eszközpozícióit rögzítik 30~perces időfelbontással, 75~napon át, egy 200$\times$200-as rácshálóra diszkretizálva. A megfigyelési terület közvetlenül nem nyilvános, ám térbeli sablonillesztéssel egyértelműen Nagoya nagyvárosi régiójához (Aichi prefektúra, Japán) köthető \cite{pinter2024revealing}.

Az adathalmaz Creative~Commons~BY 4.0 licenc alatt érhető el a Zenodo platformon\footnote{\url{https://zenodo.org/records/14219563}}. A dolgozat a \texttt{v3} adatverziót használja.

\begin{table}[H]
\centering
\caption{Az YJMob100K adatkészlet kulcsparaméterei.}
\label{tab:yjmob_summary}
\begin{tabular}{l l}
\hline
\textbf{Tulajdonság} & \textbf{Érték} \\
\hline
Forrás                  & Yahoo Japan Application, Yahoo Japan Corporation \\
Hivatkozás              & Yabe et al.\ \cite{yabe2023metropolitan} \\
Megfigyelési terület    & Nagoya, Aichi prefektúra, Japán \\
Időtartam               & 75~nap \\
Időfelbontás            & 30~perc (48~időrés/nap) \\
Térbeli felbontás       & 200$\times$200 cellás rács \\
Cellaméret (becsült)    & $\approx 455$~m (K--Ny) $\times \approx 556$~m (É--D) \\
Teljes kiterjedés       & $\approx 91$~km $\times \approx 111$~km \\
Felhasználók (dataset1) & 100~000 \\
Felhasználók (dataset2) & 25~000 (független kohorsz) \\
Rekordok (dataset1)     & 111\,535\,175 \\
Rekordok (dataset2)     &  29\,389\,749 \\
Összes rekord           & 140\,924\,924 \\
Licenc                  & CC~BY~4.0 \\
\hline
\end{tabular}
\end{table}

\subsection{Adatszerkezet}

Az alapadatfájlok (\path{yjmob100k-dataset1.csv}, \path{yjmob100k-dataset2.csv}) szerkezete azonos: minden sor egyetlen pingelési eseményt rögzít öt mezővel: \texttt{uid} (pszeudoanonim azonosító), \texttt{d} ($0\le d\le74$, megfigyelés napja), \texttt{t} ($0\le t\le47$, 30~perces időrés), \texttt{x}, \texttt{y} ($1\le x,y\le200$, rácskoordináták).

Fontos sajátosság, hogy az $x$ tengely a north--south (\textit{northing}) iránynak, az $y$ tengely az east--west (\textit{easting}) iránynak felel meg. Ez ellentétes a Descartes-koordináták szokásos $(x, y) = (\text{easting, northing})$ konvenciójával, és a feldolgozó kód minden ponton explicit módon kezeli (lásd 7.\ és 8.\ fejezet). Az adatfelvétel sem napszakot, sem dátumot nem rögzít közvetlenül; a hét napjainak ($d$~modulo~7) megfeleltetését empirikusan, az aktivitásmintázat heti periodicitásából vezettem le (lásd 9.7.\ szakasz).

\subsection{A két részadatkészlet közötti kapcsolat}
\label{subsec:two_datasets}

A Zenodo-adatközlés (v3) szerint a két fájl \emph{két diszjunkt felhasználói kohorszot} ír le ugyanazon a 75~napos időablakon és rácsterületen: a \texttt{dataset1} egy üzleti szokásrend szerinti („business-as-usual") populációt (100\,000~fő), a \texttt{dataset2} pedig egy ettől független, kisebb populációt (25\,000~fő) tartalmaz, amelynek utolsó 15~napját (\(d=60\ldots74\)) egy nem dokumentált kollektív viselkedési zavar (\enquote{vészhelyzet}) érinti.

Fontos figyelmeztetés: az \texttt{uid} mező mindkét fájlban 0-tól indul és \emph{kizárólag a saját fájljára érvényes} -- közös személyazonosítóként nem használható. A numerikus értékkészletek átfedése (\texttt{dataset2}: $0\ldots24\,999$ \(\subset\) \texttt{dataset1}: $0\ldots99\,999$) önmagában nem jelent személyazonosság-egyezést. Ezt empirikusan is ellenőriztem: a két fájl azonos \texttt{uid}-jeihez detektált otthoni cellák közül mindössze 2 esett egybe (a 23\,953 közös azonosítóból), az átlagos rácstávolságuk pedig $\approx 75$ cella ($\approx 38$~km), ami a véletlen-baseline ($\approx 104$ cella) közelében van. Ha a két fájl ugyanazokat a személyeket írná le, a home-cellák túlnyomó többségének egybe kéne esnie a tartózkodási hely napról-napra való stabilitása miatt.

Ennek két következménye van a módszertanra: (1) a vészhelyzet hatását \emph{populációszintű} (between-subject) összehasonlítással vizsgálom (\texttt{dataset1} baseline, \texttt{dataset2} kezelt; 7.4.\ és 9.5.\ szakasz); (2) az anomáliadetekciós modellt a \texttt{dataset1} normál mintáján tanítom és a \texttt{dataset2} vészhelyzeti időszakára alkalmazom, ahol a populációs heterogenitást kontrollálni kell (7.4.\ szakasz).

\subsection{Kísérő POI-fájlok}

Az adathalmaz két melléktáblát is tartalmaz: \texttt{cell\_POIcat.csv} (221\,000~sor; \texttt{(x, y, category, POI\_count)}) cellánkénti POI-kategória-eloszlással, és \texttt{POI\_datacategories.csv} (85~sor) a kategórianevek feloldásával (pl.\ \emph{Bakery}, \emph{Hospital}, \emph{Transit Station}).

A 85 finomszemcsés POI-kategóriát a \texttt{schema.py} 9 magasabb szintű \textbf{zónába} (\textit{food, shopping, leisure, transit, parking, healthcare, education, services, industry}) sorolja össze, amely táblát a \texttt{dim\_cell} dimenzió-tábla \texttt{poi\_zone} oszlopa tartalmazza. A teljes leképezés a \emph{Függelék}~A.1.\ alatt található.

\subsection{Adatminőségi és értelmezési megfontolások}

Az YJMob100K, mint minden anonimizált CDR-jellegű adathalmaz, a 3.7.\ szakaszban felsorolt általános korlátokon túl néhány specifikus megfontolást igényel: \textbf{reprezentativitás} (kizárólag a Yahoo Japan Application felhasználói, fiatalabb többséggel); \textbf{térbeli felbontás} (a névleges 500~m megtévesztő, valódi cellaméretek $\approx 454{,}5 \times 555{,}6$~m); \textbf{időbeli horgonyzás hiánya} (csak relatív heti/napi mintázatok); \textbf{esemény-jellegű hiányosság} (a vészhelyzet típusa rejtett, az anomália csak strukturálisan értelmezhető).

A korlátok ellenére az YJMob100K az egyik legnagyobb, nyíltan elérhető nagyvárosi mobilitási adatkészlet, amelynek mérete (141~M rekord) és időbeli mélysége (75~nap) a CDR-szakirodalomban szokásosnál is részletesebb elemzéseket tesz lehetővé.
